%=======================================================================================
% MERGED METHODOLOGY CHAPTER -- 2026-08-09 - 11:55pm
% Source of the section/heading layout: the dissertation.tex skeleton pasted into the planning
% session (chapter/section names preserved as given).
% Source of the facts: documentation/august_draft/4_Chapter_methodology/methodology_chapter_draft_cited.tex
% (already evidence-checked against code -- file paths, exact hyperparameters, correction dates),
% cross-checked again here against the same code files where a claim looked load-bearing.
% Still planning notes, not dissertation prose: short bullets/fragments, Purpose/Method/Evidence/
% Finding labels, NEEDS CHECKING / PENDING / SUPERSEDED flags kept exactly as found -- do not
% smooth these into narrative prose yet.
% Results-chapter numbers (accuracy tables, ablation verdicts) are deliberately NOT duplicated
% here -- every such fact below points to "Results §5.x" instead, per the source draft's own rule.
% The results ledger artifact was explicitly excluded as evidence for this file (user instruction,
% 2026-08-09 11:55pm: "the ledger hasnt been created ... work on the other sections") -- every fact below
% traces to code, `outputs/`, or `documentation/experiment_log.md`/`progress_notes.md`, not the
% ledger.
%=======================================================================================

\chapter{Methodology}

\section{Experimental questions and overall design}

Four questions actually structure every experiment in this chapter (kept as four, not three --
the pasted skeleton's three-question list omits the environment/biological-behaviour split that
the rest of this chapter and Results depend on):

\begin{enumerate}
  \setlength{\itemsep}{0pt}
  \item Which model predicts held-out top height best -- and does environmental information or
  Chapman--Richards (CR) conditioning improve that prediction?
  \item Does adding CR physics improve held-out prediction, and separately, does it improve
  predicted biological behaviour? \textbf{Two separate sub-claims, not one} -- answered by two
  different experiments with different statistical power (Experiment A vs. Experiment B below).
  Do not read one as evidence for the other.
  \item Which environmental/management conditions align with persistent, plot-level departure
  from a single shared CR curve? (Avenue 1)
  \item Which conditions align with plot-specific deviation in estimated growth trajectory, and
  does this give stronger or different attribution than Avenue 1? (Avenue 2)
\end{enumerate}

\textbf{Pipeline note.} Target rebuilt as \texttt{elev\_percentile\_95th} (raw 95th-percentile
LiDAR height, no $\times1.1$ correction); \texttt{yldc} removed from every predictive feature list
(kept only as an audit/evaluation column, and as the ingredient for Avenue 2's target -- \S4.8).
The older \texttt{Top\_Height99}/\texttt{yldc}-dependent pipeline is retired, exploratory-only
context (\S4.11).

\textbf{Splits.} Primary: \texttt{spatial\_block} (whole-compartment holdout, 60\,m train buffer).
\texttt{temporal}/\texttt{temporal\_narrow\_gap}: secondary, harder generalisation test (trains on
earlier survey years, tests on a real future survey never trained on). Stronger confirmation of
\texttt{spatial\_block}: pooled five-fold compartment-held-out spatial CV
(\texttt{spatial\_block\_kfold}), whole-population predictions concatenated across folds, not an
average of five separate R² values -- \texttt{plot\_level} (random 60/40 split, no spatial holdout
at all) is kept only as an easy-case sanity floor, not a generalisation claim.

\textbf{Terminology correction, worth stating plainly:} "environmental PINN" is not a third
avenue. It is \texttt{pinn\_env\_terrain}/\texttt{pinn\_env\_terrain\_k} -- terrain/wind-conditioned
\emph{predictive} models (\S4.2's third subsection), evaluated as accuracy, not attribution.
Avenue 1 and Avenue 2 (\S4.4 below) are the only two attribution strands.

\begin{table}[!htbp]
\centering
\scriptsize
\setlength{\tabcolsep}{3pt}
\renewcommand{\arraystretch}{0.95}
\begin{tabular}{p{2.3cm}p{1.4cm}p{3.2cm}p{4.2cm}p{2.5cm}}
\hline
Analysis strand & Analytical unit & Target & Main purpose & Models \\
\hline
Predictive comparison & Plot-year & \texttt{elev\_percentile\_95th} & Compare predictive accuracy, predict height at a given plot/survey year & CR, linear, RF, average-by-age, DNN, PINN \\
Environment-conditioned prediction & Plot-year & \texttt{elev\_percentile\_95th} & Test whether site/environmental information improves prediction or changes fitted growth behaviour & \texttt{dnn\_env\_terrain}, \texttt{pinn\_env\_terrain}, \texttt{pinn\_env\_terrain\_k} \\
Avenue 1: shared-curve attribution & Plot & \texttt{mean\_cr\_residual} & Identify conditions associated with persistent departure above/below one common growth curve & Elastic Net, XGBoost, NLME (two-stage mixed model) \\
Avenue 2: plot-curve attribution & Plot & \texttt{local\_y\_max\_difference} & Identify conditions associated with plot-specific deviation in estimated growth trajectory & Elastic Net, XGBoost, GNNWR \\
\hline
\end{tabular}
\end{table}

Outcomes (accuracy numbers, ablation verdicts, significance results) are not in this chapter --
see Chapter 5 (Results).

%------------------------------------------------
\section{Height prediction}

Both spatial and temporal generalisation are tested for every model in this section, under the
splits defined above; the spatial holdout is primary, temporal is a secondary, harder check.

\subsection*{Reference and baseline models}

\begin{itemize}
\setlength\itemsep{0pt}
\item \textbf{\texttt{average\_by\_age}}: lookup table, mean training-set height per rounded
integer age; unseen ages fall back to the overall training mean. No parametric form -- the floor
for "does any model beat a naive age lookup."
\item \textbf{Chapman--Richards curve} -- see equation and parameters below; the mechanistic/physics
floor.
\item \textbf{\texttt{linear\_baseline}}: OLS on \texttt{Age, CanopyCover, Thin,
time\_since\_thinning, time\_since\_thinning\_missing, recent\_thinning\_5yr, thinning\_status}
(one-hot, \texttt{drop\_first}, reference category \texttt{never\_thinned}). Tests whether
nonlinear capacity is needed at all.
\item \textbf{\texttt{rf\_baseline}}: \texttt{RandomForestRegressor(n\_estimators=100,
random\_state=42)}, sklearn defaults, explicitly untuned in code -- a reference point, not a
competitive model.
\end{itemize}

All baselines guarded against split-column leakage (\texttt{cpmt}/\texttt{blk}/\texttt{scpt}
never enter as features).

\textbf{Chapman--Richards functional form (as implemented, no $y_0$ term):}
\[
H(a) = y_{\max}\left(1 - e^{-ka}\right)^p
\]
Meaning of the three parameters: $y_{\max}$ -- asymptotic height; $k$ -- growth-rate constant;
$p$ -- shape parameter. Fit via \texttt{scipy.optimize.curve\_fit}, 5 random initial-value
restarts, lowest-SSE retained. Bounds: $y_{\max}\in[1.001\times\max(\text{obs}),\,5\times\max(\text{obs})]$,
$k\in[0.0001, 2.0]$, $p\in[0.01, 15.0]$.

Analytical derivative (used directly in the PINN physics loss, \S4.3):
\[
\frac{dH}{da} = y_{\max}\, p\, k\, \left(1-e^{-ka}\right)^{p-1} e^{-ka}
\]

\textbf{Anchor leakage -- found and fixed for the predictive pathway.} A pooled CR anchor (fit
once, on a 60/40 \texttt{plot\_level} random split) was originally reused across every downstream
split -- this overlaps whichever plots a given split later calls test.
\textbf{Fix:} CR refit per \texttt{split\_type}, train-partition only
(\texttt{spatial\_block}/\texttt{temporal}/\texttt{temporal\_narrow\_gap} each get their own
anchor, loaded back by \texttt{split\_type}). All PINN physics-loss anchors are now split-matched.
Measured size of the leak: Results \S5.1 (\texttt{NEEDS CHECKING} against outputs at
write-up time -- confirm the exact percentage cited there still matches the current anchor files).
A confirmatory 5-seed rerun of the no-env physics-weight comparison under the corrected anchor was
run 2026-08-09 -- outcome in Results \S5.2, not here.
\textit{Separate, still-open issue, distinct from the leak fix:} the CR fit itself pools
repeat-survey rows per plot as independent observations -- no plot-level random effect in the
curve fit.

\subsection*{DNN and PINN}

Shared base architecture (\texttt{NoEnvNetwork}, used by both \texttt{dnn\_noenv} and
\texttt{pinn\_noenv}): 3 hidden layers, 128 neurons each, LeakyReLU; Adam (lr=$10^{-4}$,
weight\_decay=$10^{-5}$); gradient clipping (max\_norm=1.0); \texttt{ReduceLROnPlateau} (factor
0.8, patience 15); early stopping on 5-epoch smoothed validation loss. \texttt{dnn\_noenv} loss:
$\text{MSE(pred, target)} + \text{L1 penalty (coef } 10^{-5}\text{)}$ -- no physics term.

\textbf{PINN loss} (real forward passes only, no synthetic collocation points):
\begin{align*}
L_{\text{data}} &= \text{MSE}(H_i, \hat H_i) \\
L_{\text{deriv}} &= \text{MSE}\!\left(\tfrac{\partial \hat H}{\partial a}\Big|_{\text{autograd}},\;
  \tfrac{dH}{da}\Big|_{\text{CR, frozen } y_{\max},k,p}\right) \\
L_{\text{traj}} &= \text{MSE}\!\left(\tfrac{\hat H_{\text{later}}-\hat H_{\text{earlier}}}{\Delta a},\;
  \tfrac{dH}{da}\Big|_{\text{CR, age}_{\text{mid}}}\right) \\
L_{\text{total}} &= L_{\text{data}} + \lambda_d L_{\text{deriv}} + \lambda_t L_{\text{traj}} + L_1
\end{align*}

Plain-language: derivative loss compares the network's own age-derivative (via autograd) against
the frozen CR curve's derivative at the same age; trajectory loss compares the network's predicted
change between two \emph{real} consecutive surveys of the same plot against the CR derivative at
the midpoint age -- built only from real survey pairs, never synthetic ones.
\textbf{Trajectory-pair leakage rule:} a pair is only used if \emph{both} endpoints are labelled
train under the active split -- verified in code to hold for spatial and temporal splits alike; no
val/test row enters a trajectory pair.
\textbf{Validation loss} is the data term only, in every PINN variant -- early stopping tracks
predictive accuracy, not physics agreement.
CR anchor inside the physics terms: split-matched (per the leakage fix above) for every final PINN
comparison; the pooled anchor was used earlier and is superseded.

\textbf{Tested physics weights, batch size, optimiser, early stopping, seeds:} see the ablation
grid in \S4.4 (Experiment A) for the actual weight values tested; batch size and optimiser as
above (shared architecture); seeds 42--46 where a seeded comparison was run (Results \S5.2 states
which comparisons currently have seed coverage and which are still pending). Full architecture and
hyperparameter table: \textbf{PENDING} -- move to appendix once finalised.

\subsection*{Environment-conditioned models}

Two structurally different mechanisms, not one "environmental PINN":

\begin{itemize}
\setlength\itemsep{0pt}
\item \textbf{Environmental DNN} (\texttt{dnn\_env\_terrain}): terrain/wind features concatenated
directly into the main network's input (same architecture as \texttt{dnn\_noenv}) -- deliberately
so the DNN sees the same information the PINN gets, for a fair "does physics help" test.
\item \textbf{$y_{\max}$-conditioned PINN} (\texttt{pinn\_env\_terrain}): main age network
unchanged; terrain/wind reach the model only through a small side network (\texttt{YMaxSubNetwork}:
1 hidden layer, 16 units, LeakyReLU) producing an \emph{additive} adjustment to the frozen anchor:
$y_{\max,\text{row}} = y_{\max,\text{global}} + \Delta_{\text{subnet}}(\text{terrain})$. Design
rationale cited in code: environment sets the growth \emph{ceiling}, not the trajectory shape
(Socha et al. 2021, ADA/GADA).
\item \textbf{$y_{\max}+k$-conditioned PINN} (\texttt{pinn\_env\_terrain\_k}): adds a second,
separate sub-network conditioning $k$ \emph{multiplicatively in log-space} to keep it positive:
$k_{\text{row}} = k_{\text{global}}\cdot \exp(\Delta_{\text{subnet}}(\text{terrain}))$. $p$ stays
global/frozen in every variant (metabolic-scaling-theory justification in code). Code explicitly
flags that the two sub-networks could become confounded (both driven by the same terrain inputs)
-- the \texttt{freeze\_y\_max} ablation tests this directly (finding: Results \S5.3.3).
\end{itemize}

Clarification per the original plan's own instruction: these three models are still
\textbf{prediction} models -- they are not the attribution models (\S4.4).

\textbf{Batch size, a resolved confound.} \texttt{dnn\_env\_terrain.py}'s file-level default is
512 (kept equal to \texttt{dnn\_noenv.py}'s own default, for isolating the terrain-feature effect
within the DNN family), but every actual fit behind the reported
\texttt{dnn\_env\_terrain}-vs-\texttt{pinn\_env\_terrain}/\texttt{pinn\_env\_terrain\_k}
comparisons explicitly overrides this to \texttt{--batch-size 256}, confirmed matched directly
against \texttt{outputs/run\_logs/} -- \textbf{not} an open confound (this corrects an earlier,
now-superseded framing of this as unresolved).

Architecture diagrams (3, to be drawn -- described precisely enough in the source draft that a
figure can be produced from text alone): (1) \texttt{dnn\_noenv}/\texttt{pinn\_noenv} single-box
network with two dashed physics-loss-only branches (derivative, trajectory); (2)
\texttt{dnn\_env\_terrain} -- same box, wider input layer, no physics branches; (3)
\texttt{pinn\_env\_terrain}/\texttt{pinn\_env\_terrain\_k} -- main box architecturally identical to
(1) with \textbf{no} arrow from terrain/wind into it; terrain/wind instead feed the small
$y_{\max}$(/$k$) side network(s), which feed only the dashed CR-reference-curve box. \textbf{PENDING}
-- not yet drawn.

%------------------------------------------------
\section{Environmental Attribution methods}

"Attribution" here means held-out statistical association between environmental/management
predictors and a derived growth-deviation target -- \textbf{not} causal evidence. Its practical
purpose in this project is also a screening question: whether a variable's resolution/spatial
coverage/signal is good enough to be worth including at all (\S4.5 below is the actual screening
method used; distinct in Avenue 1 vs. Avenue 2).

\textbf{Predictor groups, scaling/encoding, leakage removal -- shared across both avenues:}
\begin{itemize}
\setlength\itemsep{0pt}
\item \texttt{Age} excluded from Avenue 1 -- circular, it is the CR curve's only input and
\texttt{mean\_cr\_residual}'s sole construction variable.
\item Neighbour-height variables (\texttt{neighbour\_mean\_height},
\texttt{neighbour\_height\_differential}) removed after confirming they leak test-set ground truth
(built from other plots' real heights within 75\,m, pre-split; 95.9\% of a test plot's neighbours
are themselves test-set plots). Measured effect of removing them: Results \S5.1.
\item Management kept as its own category throughout
(\texttt{stand\_structure}: \texttt{CanopyCover, Thin, time\_since\_thinning,
time\_since\_thinning\_missing, recent\_thinning\_5yr}), never folded into "environmental."
\end{itemize}

\textbf{Design principle: a variable's route into a model determines how it is screened.} Three
purpose-built selection procedures are used, not one uniform rule, because a redundant input
costs a model different amounts depending on what happens to that input downstream. A neural
network's hidden layers can absorb two correlated raw inputs without their collinearity
corrupting anything the network reports; a linear coefficient or a SHAP value cannot -- collinear
inputs there directly distort the reported attribution. Each of the three procedures below is
scoped to exactly the model family it feeds, on that basis.

\textbf{(i) Avenue 1's predictive feature tiers} (\texttt{ENV\_TERRAIN\_FEATURE\_SETS} in
\texttt{models/common/torch\_data.py}) feed \texttt{dnn\_env\_terrain}/\texttt{pinn\_env\_terrain}/
\texttt{pinn\_env\_terrain\_k} directly, and are built as cumulative stages by design --
terrain-only ($\texttt{stage1\_terrain}$, 13 vars incl. multiscale TPI) $\to$ +wind
($\texttt{stage2\_terrain\_wind}$, +8 vars) $\to$ +whichever of climate/soil/edge-position carries
real individual signal ($\texttt{stage3\_terrain\_wind\_plus}$) $\to$ +every remaining category
unfiltered ($\texttt{stage4\_all\_environmental}$), so that a model comparison across tiers
answers "how much does each successive category of information add," not "which single variable
matters." Deliberately \textbf{no VIF step at this stage}: per the point above, redundant raw
inputs are not a cost this tier needs to pay to control. The one gate applied before a category
enters stage 3 is a per-category minimum-signal check, not a strength ranking: that category's
single strongest column must clear $|\text{Spearman }\rho|\ge 0.10$ against the target
(\texttt{STAGE3\_SIGNAL\_THRESHOLD}, \texttt{multicollinearity\_screen\_av1.ipynb}) -- set
deliberately low, to exclude only a category with no individual signal at all, not to rank
categories by strength. \texttt{stage3\_terrain\_wind\_plus} and \texttt{stage4\_all\_environmental}
resolve to the identical variable list under this run's data, since every remaining category
(climate, soil/site, edge-position) cleared that gate -- kept as two separately named tiers
regardless, since the gate is a real check applied fresh every time the screen is re-run, not a
result guaranteed to repeat.
(Legacy tiers \texttt{terrain\_wind\_solid} (5 vars, the 2026-07-31 working default),
\texttt{terrain\_wind\_extended}, \texttt{broad} and \texttt{broad\_legitimate} predate this
staged design and are retained in code for continuity/comparison, not as the current tier
system.)

\textbf{(ii) Avenue 1's attribution screen} (\texttt{multicollinearity\_screen.py}, in
\texttt{models/xgb\_environmental/}) feeds the models whose coefficients or SHAP values are
reported directly -- Elastic Net, XGBoost and NLME, fitting \texttt{mean\_cr\_residual} -- and is
the properly collinearity-controlled companion to (i), used exactly where collinearity control
actually matters. Three stages: (1) drop deterministic duplicates -- variables that are an exact,
documented closed-form function of another candidate already in the pool (e.g.
\texttt{inverse\_slope\_proxy = -slope\_degrees}), which by definition carry zero new information;
(2) drop near-exact empirical duplicates, $|\text{Spearman }\rho|\ge 0.95$, keeping one
representative per near-duplicate pair -- tie-break prefers an externally-measured variable over a
derived one, then whichever has the higher $|\rho|$ with the target, computed fresh; (3) iterative
VIF reduction, threshold 5.0 (Dormann et al. 2013) -- drop the single worst-VIF column, recompute
every remaining column's VIF from scratch, repeat, so that removing one variable's effect on every
other variable's score is never missed (capped at 50 iterations). A fourth, \textbf{non-dropping}
pass then compares each surviving variable's VIF computed on the raw data against the same VIF
computed after subtracting each compartment's own mean: a variable that is high-VIF raw but
comfortably low-VIF once that shared regional trend is removed is tracking the same broad spatial
pattern as its neighbours (e.g. elevation and temperature, via the lapse rate), not literally
duplicating a measurement -- flagged for direct inspection rather than auto-dropped, since that is
a real, defensible variable, not a duplicate.

\textbf{(iii) Avenue 2's scope system} (\texttt{FEATURE\_GROUPS}/\texttt{SCOPE\_GROUPS} in
\texttt{models/growth\_curve\_attribution/broad\_environmental\_check.py}) feeds Elastic Net,
XGBoost and GNNWR, fitting \texttt{local\_y\_max\_difference}, and is built from the group up
rather than the variable up: five named groups (\texttt{terrain\_wind} 17 vars,
\texttt{climate} 5, \texttt{soil\_site} 5, \texttt{edge\_position} 6, \texttt{management} 5) are
each decided once, then combined directly into the tested scopes
(\texttt{terrain\_wind}; \texttt{terrain\_wind\_plus\_management}; \texttt{broad\_environment} =
terrain\_wind+climate+soil\_site+edge\_position, deliberately excluding management so its own
contribution stays visible rather than folded silently into "environmental";
\texttt{broad\_environment\_plus\_management} = every category combined). This is the right level
for the question Avenue 2 asks of these scopes -- how much does adding a whole category help --
so no further collinearity pass runs across categories at this stage; the finer question of
whether individual variables within a fitted scope behave stably is answered afterward, by
gain-based rank stability across folds and by SHAP (\S4.6/\S4.9), not by pre-filtering here.
What \emph{is} decided in advance, once, is each group's own internal representation: candidate
variants of the same underlying signal (e.g. Global Wind Atlas wind speed vs. local shelter
terms \texttt{topex}/\texttt{windward\_topex}/\texttt{whcl}; TPI at 250\,m/500\,m vs. local
relief) are each swapped into a fixed base set, Elastic Net and XGBoost refit, and the
representation with the larger held-out R² gain over baseline is the one that goes into
\texttt{terrain\_wind} (\texttt{correlation\_screen.py} reports Spearman among the candidates
first, as a cheap check, but the refit-and-compare result is the actual decision rule --
\texttt{feature\_representation\_check.py}). This representation-choice step is model-fit-based
by design, deliberately different from (ii)'s fixed VIF cutoff: it answers "which version of
this signal predicts best," a question a correlation threshold alone cannot.

\textbf{SHAP is not a selection step in either avenue.} It is applied once, after fitting is
complete, purely to interpret an already-chosen model on an already-decided scope
(\texttt{explain\_signal.py}: XGBoost refit on \texttt{local\_y\_max\_difference} over the full
population, no held-out split -- an interpretation-only fit, \S4.6/\S4.9) -- never to decide which
variables enter a tier or scope in the first place.

\textbf{Why Avenue 1 has no separate management tier and Avenue 2 does, by design, not
inconsistency:} Avenue 1's stand/management variables are baseline inputs on every Avenue 1
model -- fed directly into the DNN/PINN main network for the predictive tiers, and always present
in the attribution feature list -- so there is nothing to isolate as an optional addition.
Avenue 2's attribution models carry no such fixed baseline, so management is deliberately kept
isolatable as its own category, to make its contribution visible rather than silently absorbed
into "environmental" (\texttt{broad\_environment} vs.\ \texttt{broad\_environment\_plus\_management}
is exactly this comparison).

Which representations won each check, and how the tiers/scopes performed against each other:
Results \S5.5/5.6, not here.

\subsection{avenue 1: shared curve residual}

\textbf{Target:}
\[
\bar r_i = \frac{1}{T_i}\sum_t \left[H_{it} - H_{\text{CR}}(a_{it})\right]
\]
(\texttt{mean\_cr\_residual}) -- whether a plot was, on average, taller or shorter than the shared
curve predicts. CR anchor here is the \textbf{pooled} anchor (fit once, on the
\texttt{plot\_level}-split train partition, $\approx$60\% of rows), applied uniformly to every
plot-year regardless of downstream split -- a deliberate design choice, distinct from the
split-matched anchor used in the PINN physics loss (\S4.2): \texttt{mean\_cr\_residual} is the
fixed shared yardstick that \emph{defines} the attribution target, not a model input, so a single
consistent reference curve is the point, not a leakage risk in the same sense.

\textbf{Models:}
\begin{itemize}
\setlength\itemsep{0pt}
\item \textbf{Elastic Net} -- regularised linear associations.
\item \textbf{XGBoost} -- nonlinear relationships and interactions.
\item \textbf{NLME} (mixed-effects) -- a two-stage \emph{approximation}, not a true nonlinear
mixed-effects fit. Stage 1: fixed pooled CR curve. Stage 2: \textbf{linear} mixed model
(\texttt{statsmodels MixedLM}) on the resulting residual, standardised fixed effects,
\textbf{random intercept per compartment only} (no random slope). Coefficients reported under
REML; variance-explained comparisons (null vs. full model) deliberately refit under ML, since REML
likelihoods are not comparable across differing fixed-effects structures (Pinheiro \& Bates). An
earlier NLME version used fewer variables, included a variable later dropped by VIF screening, and
an uncorrected REML variance comparison -- \textbf{SUPERSEDED} by the current 18-variable
VIF-screened rebuild (both versions' numbers: Results \S5.5).
\end{itemize}

\textbf{Evaluation:} five-fold compartment-held-out spatial CV, same 60\,m buffer, separate
validation fold. Feature-set naming correction: the maintained comparison is
\texttt{terrain\_and\_wind\_only} vs. \texttt{all\_environmental} -- \texttt{all\_environmental}
already includes the management/stand-structure variables; there is no separate
"all\_environmental\_plus\_management" tier in code.

\textbf{Statistical confirmation:} paired fold-level test (Wilcoxon signed-rank + paired t-test on
the 5 fold-level R² values). Weaker than a compartment-independent cluster bootstrap -- uses only
the fold-level R² summaries already on disk, not raw per-compartment predictions, because no
fitted model artefacts were saved and retraining was out of scope for this check. \textbf{No
compartment-cluster bootstrap currently exists for Avenue 1's EN/XGBoost/NLME comparisons} --
flagged as the strongest concrete next addition (Avenue 2 has the fuller treatment, \S4.5 below).

\textbf{Spatial autocorrelation of the residual:} global Moran's I via \texttt{esda}/\texttt{libpysal},
\texttt{DistanceBand} weights with threshold set from the fitted semivariogram range (not guessed),
999 permutations. Local Moran's I (LISA) via KNN weights ($k$=20 and $k$=50 as a sensitivity pair),
Benjamini--Hochberg FDR correction on the permutation p-values (not a flat $p<0.05$ per plot).
Result, and comparison to the other five models: Results \S5.2.1/\S5.5.

\subsection{avenue 2: plot curve deviation}

\textbf{Target construction (corrected detail vs. the original plan sketch).} Per-plot $y_{\max}$
is fit with shape parameters $(p_4, p_5)$ held fixed per plot; because
$H = y_{\max}(1-e^{-p_4 a})^{p_5}$ is linear in $y_{\max}$ given fixed shape, the fit is
\textbf{closed-form} least-squares regression through the origin
($y_{\max,\text{fit}} = \sum(H\cdot\text{shape\_term}) / \sum(\text{shape\_term}^2)$), not
iterative \texttt{curve\_fit} -- and, unlike Avenue 1's pooled anchor, uses only that plot's own
rows, so there is no cross-plot self-influence pathway analogous to \S4.4's Avenue 1 target.

$y_{\max}$ implied by yield class is \textbf{derived analytically}, not looked up: rearranging the
Forestry Commission's own \texttt{GYCspec95} formula,
$y_{\max,\text{yldc}} = \text{yldc}\cdot p_2 + p_1 + p_3\cdot 2$, using each plot's own recorded
$(p_1..p_5)$ tuple (7 distinct value-sets in the data), not one universal species constant.

\textbf{Target:} $\Delta y_{\max,i} = y_{\max,i}^{\text{local}} - y_{\max,i}^{\text{expected}}$
(\texttt{local\_y\_max\_difference}). CR is used only to build the target -- never inside the
attribution-model loss itself.

\textbf{Local curve fitting, eligibility, uncertainty:}
\begin{itemize}
\setlength\itemsep{0pt}
\item Disturbance-cleaning rules (interval-level, on consecutive-survey drops in
height/canopy/volume): \texttt{height\_collapse} (height\_drop\_frac $\ge 0.25$);
\texttt{clearfell\_like} (\texttt{height\_collapse} \& canopy\_drop\_frac $\ge 0.70$ \&
volume\_drop\_frac $\ge 0.70$); \texttt{measurement\_inconsistent} (\texttt{height\_collapse} \&
canopy\_drop\_frac $\le 0.10$ \& volume\_drop\_frac $\le 0$); \texttt{ambiguous\_disturbance}
(\texttt{height\_collapse} but neither of the above -- \textbf{retained} in the modelling
population, could reflect a real environmental mechanism, e.g. partial windthrow, not pure noise).
\item Exclusion rule: a plot with \emph{any} \texttt{clearfell\_like} or
\texttt{measurement\_inconsistent} interval is dropped entirely from curve fitting (the fit uses
all of a plot's rows jointly).
\item Fitting-error/uncertainty from only 4 or 6 observations per plot: \textbf{PENDING} --
target-reliability figures (exclusion counts, distribution of local $y_{\max}$, boundary/implausible
estimate frequency) not yet consolidated anywhere in current artefacts (Results \S5.6.1 flags this
as an open gap, not fabricated).
\end{itemize}

State clearly, per the source draft's own framing: Elastic Net and XGBoost are fitted as
\textbf{global} comparison models for Avenue 2 -- same model descriptions as Avenue 1 above, not
repeated.

\textbf{GNNWR} (Geographically Neural Network Weighted Regression, Du et al. 2020; third-party
\texttt{gnnwr} package, not reimplemented):
\[
\hat y_i = \beta_0(s_i) + \sum_j \beta_j(s_i)\, x_{ij}
\]
GNNWR allows the predictor--target relationship to vary with location $s_i$; Elastic Net/XGBoost
estimate one global relationship regardless of location. Mechanism: a spatial-weighting sub-network
takes each plot's distance-to-reference-point vector and outputs a plot-specific coefficient vector,
which then combines linearly with the plot's own feature vector.

Current evaluation trains on the \textbf{full} training population via pooled five-fold spatial CV,
the same fold assignment and pooling convention as EN/XGBoost -- a fully apples-to-apples
comparison. An earlier phase capped the reference set via compartment-stratified subsampling as a
GPU-memory workaround -- \textbf{SUPERSEDED} by the full-population run (both versions' numbers:
Results \S5.6). A patched diagnostics routine skips the O($n^2$) "hat matrix" above 2,000 rows;
this affects only AIC/AICc/F-test-style diagnostics, not the reported R²/RMSE, which are computed
independently of the hat matrix.

\textbf{Evaluation:} five-fold compartment-held-out spatial CV, same 60\,m buffer, across static
scopes \texttt{terrain\_wind}, \texttt{broad\_environment}, \texttt{terrain\_wind\_plus\_management},
\texttt{broad\_environment\_plus\_management} (\texttt{broad\_environmental\_check.py}; groups:
terrain\_wind 17 vars, climate 5, soil\_site 5, edge\_position 6, management 5).

\textbf{Broad-environmental scope: SUPERSEDED, do not cite numerically.} CSVs generated 2026-08-04,
before an XGBoost evaluation-set-leak fix (commit \texttt{3ffc5fc}) landed the same day. The
notebook itself has not been regenerated post-fix. Only defensible claim until regenerated:
qualitative -- management added real signal, broader environmental variables beyond terrain/wind
added comparatively little -- not the specific R² values.

\textbf{\texttt{env\_deviation}} (a smaller, earlier module): a predecessor/sibling to this
pipeline, not formally retired -- predicts row-level residual of either a split-matched-anchor CR
curve or \texttt{dnn\_noenv}'s own prediction, using terrain-only XGBoost, no physics loss, no
joint training. Still referenced in current code (\texttt{phase0\_checks.py}, \texttt{saving.py});
described here as a parallel earlier approach, not folded into or superseded by the per-plot
\texttt{local\_y\_max\_difference} pipeline.

\textbf{Full feature list (all tiers, both avenues): PENDING -- appendix.}

%------------------------------------------------
\section{Evaluation and Interpretation}

\begin{itemize}
\setlength\itemsep{0pt}
\item Held-out MAE, RMSE, R² -- primary metrics throughout, every model family.
\item Spatial cross-validation -- five-fold compartment-held-out, 60\,m buffer, as above; single
\texttt{spatial\_block} split kept as the primary headline comparison, five-fold as the stronger
confirmation.
\item Moran's I and semivariograms for residual spatial structure -- \texttt{esda}/\texttt{libpysal}
throughout; Avenue 1 uses \texttt{DistanceBand} weights (threshold from the fitted semivariogram
range); Avenue 2 and the predictive-model residual comparison use KNN weights ($k$=8 for Avenue 2,
chosen because a fixed distance band gives wildly uneven neighbour counts, 681--7,141, in this
data).
\item Monotonicity and physics-loss diagnostics for PINNs -- proportion of trajectories remaining
monotonic, implausible negative height changes, smoothness of predicted trajectories, derivative/
trajectory residuals, predictions outside a plausible height range. \textbf{Status: mostly
unresolved as of the current results draft} -- exists as a category in this evaluation plan, but
only one specific cross-model finding exists to date (the \texttt{pinn\_env\_terrain\_k}
$y_{\max}$/$k$ sign-flip, \S4.2/\S4.3 above; Results \S5.3.3). State this plainly rather than
implying the category is fully covered.
\item SHAP, coefficients and local GNNWR coefficients for interpretation -- SHAP is Avenue 2 only
(\texttt{explain\_signal.py}: XGBoost, \texttt{n\_estimators=500, max\_depth=4, learning\_rate=0.04},
fit on \texttt{local\_y\_max\_difference} over the \textbf{full population, no held-out split} --
explicitly interpretation-only, not a performance evaluation -- then \texttt{shap.TreeExplainer},
mean $|\text{SHAP}|$ per feature, 17 features; a plain Spearman-correlation pass runs first as a
cheap sanity check). Avenue 1 reports NLME standardised coefficients and Elastic Net coefficients
directly, no SHAP. XGBoost feature-importance rank stability (gain-based, not SHAP, for cost) is
computed within each of the 5 spatial folds for Avenue 2 -- mean/std/min/max rank per variable.
\item Seed or cohort sensitivity where completed -- seed sweeps exist for the no-env physics-weight
ablation (5 seeds, 42--46); cohort sensitivity is structural throughout (4survey primary, 6survey
secondary, consistently smaller and noisier -- 47--296 compartments depending on the comparison,
treat every 6survey number as a secondary check, not equal-weight evidence).
\item Comparisons treated as controlled \textbf{only} when target, split, features and settings are
matched -- stated explicitly per comparison in Results, not assumed globally.
\end{itemize}

\textbf{Statistical significance -- what actually exists, and what is missing, stated once here so
it is not repeated per-avenue:}
\begin{itemize}
\setlength\itemsep{0pt}
\item Avenue 1: paired fold-level test only (Wilcoxon + paired t, $n=5$ folds) for
\texttt{terrain\_and\_wind\_only} vs. \texttt{all\_environmental}. No compartment-cluster
bootstrap yet.
\item Avenue 2: the fuller treatment -- paired fold-level tests \emph{and} a compartment-cluster
bootstrap ($n=2{,}000$ resamples, resampling whole compartments not rows) for the GNNWR-vs-baseline
R² difference, applied only to that one comparison. Both fold-level tests are underpowered at
$n=5$: Wilcoxon's smallest possible p-value at $n=5$ is 0.0625 regardless of effect size; the
paired t-test's normality assumption on only 5 differences is not well justified either -- both
read as descriptive of direction, not strict significance claims. The identified fix for the
$n=5$ power problem is \textbf{repeated-seed CV} (more fold-assignment seeds, not more folds --
the compartment pool does not support many more folds without shrinking each holdout too far);
not yet run, flagged as the concrete next step if stronger evidence is needed.
\end{itemize}

\textbf{Methodological limitations, kept here rather than a separate section (per the original
plan's own instruction):}
\begin{itemize}
\setlength\itemsep{0pt}
\item Only four or six observations per plot -- limits estimation of detailed plot-specific
trajectories, directly bears on Avenue 2's target reliability.
\item Uncertainty in plot-level curve fitting (\S4.4, Avenue 2) -- not yet quantified.
\item Buffer-adequacy limitation: 60\,m is set from plot spacing ($\approx3\times$ the 20\,m grid
cell), not from a measured autocorrelation range -- which is measured to be far larger (CR
residual's own semivariogram range $\approx$3,956--4,284\,m; Results \S5.1). 60\,m protects
against sub-canopy pixel-level leakage; it is not, by itself, sufficient to rule out longer-range
spatial leakage. Whole-compartment holdout, not the buffer, is what does that work -- the buffer is
a secondary control.
\item Coarse or static environmental predictors -- see Data chapter's own limitations table; not
duplicated here.
\item Attribution represents association, not causation, throughout both avenues.
\item Three claims kept explicitly separate, never treated as interchangeable evidence: (1)
implementation ran as intended; (2) physics/attribution changed biological behaviour or residual
structure; (3) held-out prediction improved. Research Question 2 (\S4.1) is the concrete
illustration of why this separation matters.
\end{itemize}

%------------------------------------------------
\section*{Approaches tried and later superseded (procedural history, not results)}

\begin{itemize}
\setlength\itemsep{0pt}
\item Retired \texttt{Top\_Height99} + \texttt{yldc} predictive pipeline -- historical context only.
\item Early DNN/PINN runs with mismatched batch size, before the 2026-07-30 rebuild.
\item Pooled-anchor PINN physics-loss comparisons, before split-matched freezing (2026-08-01).
\item Avenue 2 broad-environmental CSVs generated 2026-08-04, before the same-day evaluation-leak
fix; notebook not yet regenerated.
\item Earlier, smaller NLME model -- superseded by the 18-variable VIF-screened rebuild (2026-08-08).
\item E7 single-split 6survey R² values -- superseded by E9's five-fold confirmation.
\item E10 feature-tier sweep -- dropped from the dissertation plan entirely.
\item \texttt{temporal\_narrow\_gap} numeric results from the retired \texttt{Top\_Height99}+\texttt{yldc}
pipeline (2026-07-13 to 2026-07-20) -- not comparable to the current pipeline's limited,
smoke-tested runs.
\end{itemize}

\textbf{Open, not stale} -- unanswered rather than superseded, should not be silently dropped from
the write-up:
\begin{itemize}
\setlength\itemsep{0pt}
\item Compartment-cluster bootstrap for Avenue 1's Elastic Net/XGBoost/NLME comparisons.
\item Repeated-seed spatial CV for Avenue 2's GNNWR-vs-baseline significance test.
\end{itemize}

%=======================================================================================
